\section{Consolidation Question Q3: Compatibility with FedAvg-Style Local Training}
\label{sec:consolidation-q3}

The second experimental consolidation task tests whether V1 is a generic communication layer or is fundamentally tied to the asynchronous FedSGD-like setting used in most of Part~I. The decisive change is to replace the one-step stochastic-gradient input by a conventional local model delta generated after several client-side optimization steps.

\subsection{Synchronous FedAvg starting point}

The first study should use conventional synchronous FedAvg to isolate the effect of local training from asynchrony. At global round $r$, the server provides the current model $w^r$ to the selected clients. Client $i$ initializes
\begin{equation}
    w_{i,0}^r=w^r
\end{equation}
and performs $E$ local SGD steps,
\begin{equation}
    w_{i,e+1}^r
    =w_{i,e}^r-\eta\nabla\ell_i(w_{i,e}^r;\xi_{i,e}),
    \qquad e=0,\ldots,E-1.
\end{equation}
The conventional local FedAvg update is
\begin{equation}
    \Delta_i^r=w_{i,E}^r-w^r.
\end{equation}
For Event-FedAvg, the frozen stateful encoder receives
\begin{equation}
    \boxed{u_i^r=\Delta_i^r}
\end{equation}
rather than a stochastic gradient. The same persistent state, threshold, address/sign event payload, reset rule, and scheduled server jump are then applied.

This formulation directly tests the intended abstraction from Section~\ref{sec:algorithm-consolidation}: the compressor should operate on a local learning update, not on one special optimizer quantity.

\subsection{Local-computation sweep}

The local-work variable should be swept over
\begin{equation}
    E\in\{1,5,10\}
\end{equation}
with an optional $E=20$ stress test if the first three settings are stable. The $E=1$ case connects back to the existing FedSGD-like regime, while larger $E$ values introduce genuine client drift and test whether persistent event memory remains meaningful when successive deltas are generated around different local trajectories.

The primary hypotheses are:
\begin{enumerate}
    \item Event-FedAvg remains trainable and communication-efficient for $E>1$.
    \item Increasing local work does not cause uncontrolled event-packet growth.
    \item Persistent event state does not accumulate stale sign information so strongly that local model drift destroys convergence.
\end{enumerate}

\subsection{Baselines and benchmark sequence}

The initial benchmark should be the ten-class Fashion-MNIST MLP because it is already validated, inexpensive, and easy to instrument. The minimum comparator set is
\begin{equation}
\boxed{
\text{dense FedAvg},
\quad
\text{FedAvg + EF-TopK},
\quad
\text{FedAvg + sign/EF},
\quad
\text{Event-FedAvg}.
}
\end{equation}
The local optimizer, client sampling, number of local steps, batch size, and learning-rate tuning budget must be matched across communication methods. If the MLP study passes, one selected operating point should be repeated with the compact CNN from Experiment~15A. A new dataset is not required.

\subsection{Mandatory diagnostics}

Local training introduces failure modes that were largely absent in the one-gradient-completion setting. The experiment should record
\begin{itemize}
    \item $\norm{\Delta_i^r}$ and its distribution across clients and rounds,
    \item coordinate events and nonempty packets per client per round,
    \item events per transmitted local delta,
    \item sign reversals of emitted coordinates across rounds,
    \item residual membrane norm $\norm{z_i^r}$,
    \item client-level communication imbalance,
    \item test loss and accuracy versus both rounds and transmitted bits.
\end{itemize}
These diagnostics should establish whether failures, if any, are caused by local client drift, update-scale mismatch, or persistent residual state.

\subsection{Decision gates}

A minimum pass requires that Event-FedAvg at $E=5$ reaches essentially the same predictive regime as dense FedAvg while preserving a clear communication advantage over the closest compressed FedAvg baseline. A stronger pass requires that this conclusion also holds for $E=10$ and transfers to one CNN operating point.

The main outcomes are therefore
\begin{description}
    \item[Strong pass.] The stateful event encoder works across multiple local-step counts and one additional architecture. V1 can be described as a generic FL communication layer.
    \item[Qualified pass.] The method works for modest local training but degrades for large $E$. The algorithm remains viable with an explicit local-work regime or requires a scale normalization that must be justified separately.
    \item[Fail.] The mechanism only works when the input is approximately an instantaneous stochastic gradient. The contribution should then be framed specifically as an asynchronous FedSGD communication mechanism rather than as a generic FL compressor.
\end{description}

This question should be resolved before engineering the final bidirectional protocol, because the definition of the client update source determines what the complete FL algorithm actually communicates.
